\section*{Introdução}
\addcontentsline{toc}{section}{Introdução}
\label{sec:introducao}
% Artigo 2: H2 e H2-a. Texto adaptado do capítulo 1 original, preservado em backup.

O estoque ambiental afetado pelas decisões de uso da terra não é alheio ao processo produtivo. A vegetação nativa participa das condições biofísicas que sustentam a produção e, por essa via, pode afetar os resultados econômicos. A relação entre cobertura vegetal, evapotranspiração e precipitação encontra respaldo em evidências sobre o transporte de umidade em regiões tropicais \cite{SpracklenEtAl2012}; diferenças entre tipos de vegetação também estão associadas às propriedades físicas, químicas e microbianas do solo, como documentado no Himalaia central \cite{JoshiGarkoti2023}. Esses resultados ajudam a fundamentar mecanismos possíveis, embora sua magnitude não possa ser transferida diretamente para a agricultura mato-grossense.

Isso ocorre, sobretudo, porque o estoque local de vegetação pode modificar as condições ecológicas do território e, consequentemente, a resposta da agricultura a choques climáticos. Em paisagens agrícolas inglesas, a composição e a configuração da paisagem estiveram associadas a diferentes dimensões da resiliência do rendimento do trigo, sem que as áreas de maior rendimento relativo fossem necessariamente as mais estáveis ou resistentes a eventos extremos \cite{RedheadEtAl2020}. A questão deixa, portanto, de ser apenas se a vegetação está correlacionada com maior ou menor produtividade e passa a ser se ela altera a sensibilidade da produção a perturbações externas.

Essa distinção é particularmente importante porque variações climáticas de curto prazo podem fornecer uma fonte de variação plausivelmente exógena para investigar a resposta da atividade econômica ao clima, sob condições de identificação adequadas \cite{DellEtAl2014}. Se áreas com diferentes estoques prévios de vegetação respondem de maneira distinta a choques comparáveis de precipitação, então cabe investigar em que medida essa diferença se relaciona ao estado ambiental anterior. A anterioridade da vegetação, contudo, não elimina sua associação com solo, manejo, irrigação e escolhas passadas de localização da produção. Desta feita, a interpretação causal da interação entre clima e vegetação exige mais do que a inclusão de uma variável ambiental defasada.

Não menos importante, essa relação pode ser espacialmente heterogênea. Evidências internacionais indicam que a sensibilidade dos rendimentos ao clima depende de textura do solo, carbono orgânico e disponibilidade de água \cite{HuangEtAl2021}. Em outra escala de análise, a resistência de ecossistemas à seca também varia com cobertura e manejo, embora índices de vegetação obtidos por sensoriamento remoto não se confundam com o rendimento de uma cultura agrícola \cite{XiaoEtAl2023}. Portanto, não se pressupõe que uma mesma proporção de vegetação produza efeitos equivalentes em territórios ambientalmente distintos.

Nesse contexto, Mato Grosso constitui um caso empírico especialmente informativo para a investigação dessas relações. O estado combina produção agropecuária em larga escala com três biomas --- Amazônia, Cerrado e Pantanal --- e diferenças relevantes de vegetação, clima, disponibilidade hídrica e sistemas produtivos. Estudos sobre a intensificação agrícola e o mapeamento do uso da terra documentam essa diversidade produtiva e territorial \cite{VanWeyEtAl2013,SimoesEtAl2020}. Esses biomas representam, portanto, não apenas formações ecológicas distintas, mas também contextos nos quais pode variar a contribuição do estoque de vegetação para as condições produtivas.

Sendo assim, a heterogeneidade territorial permite investigar se o estoque de vegetação modifica a resposta da produtividade aos choques climáticos e se essa relação varia entre biomas. O objetivo não é apenas verificar se a vegetação está associada ao nível médio de produtividade, mas também se ela altera a intensidade com que condições climáticas adversas afetam os resultados produtivos. A comparação estará, por óbvio, condicionada à existência de produção e de observações comparáveis em cada recorte. A presença dos três biomas no estado não assegura, por si só, suporte para estimar três respostas agrícolas com a mesma precisão.

Na organização da tese, este artigo corresponde à hipótese H2. O capítulo 1 examina, por meio de H1, como retornos econômicos e instituições condicionam decisões de conversão ou conservação. Aqui, o foco recai sobre a contribuição produtiva do estoque ambiental resultante dessas decisões. Os dois mecanismos compõem uma relação de interdependência: decisões econômicas alteram o estado ambiental, enquanto esse estado pode alterar os resultados econômicos que fundamentam decisões subsequentes. A articulação conceitual entre os artigos não equivale, nesta etapa, à estimação de um sistema simultâneo nem a um teste das expectativas intertemporais dos agentes.

\medskip
O \textbf{problema} do artigo pode ser sintetizado no seguinte questionamento:

\noindent\emph{Em que medida o estoque de vegetação nativa previamente existente no território modifica a resposta da produtividade agrícola a choques climáticos, e como essa relação varia entre os biomas de Mato Grosso?}

\medskip
O \textbf{objetivo geral} deste artigo é investigar se o estoque de vegetação previamente existente modifica, em magnitude economicamente relevante, a resposta da produtividade agrícola a choques climáticos em contextos territorialmente heterogêneos. De forma específica, o artigo busca: (i) caracterizar a distribuição do rendimento agrícola, do estoque vegetal e da exposição a choques de precipitação; (ii) estimar a interação entre vegetação anterior à janela agrícola e choques de precipitação sobre o rendimento da soja; (iii) avaliar a heterogeneidade dessa relação entre biomas e municípios mistos, conforme o suporte dos dados; e (iv) examinar a sensibilidade das estimativas às definições de clima, vegetação, unidade territorial e condições produtivas locais.

A relevância deste trabalho decorre, em termos teóricos, de tratar o \textbf{meio ambiente} como componente do processo econômico, capaz de afetar resultados produtivos e, simultaneamente, ser transformado por decisões econômicas. A abordagem de capital natural permite relacionar a condição ambiental aos fluxos futuros de renda da agricultura, como exemplificado pela valoração dos serviços de suporte do solo \cite{Brady2015}. Essa relação envolve também uma \textbf{dimensão territorial}, na medida em que condições ambientais distintas podem produzir respostas econômicas diferentes, e uma \textbf{dimensão intertemporal}, pois decisões presentes de uso da terra alteram estoques que podem condicionar oportunidades e riscos produtivos futuros.

Em termos empíricos e aplicados, compreender essas relações pode: (a) orientar decisões de produtores sobre conservação, manejo e adaptação a choques climáticos; (b) esclarecer uma das fontes de variação e de risco do rendimento agrícola; e (c) subsidiar a consideração das condições locais no desenho de políticas de conservação. Tais contribuições não dependem de presumir a ineficácia das regras ambientais ou dos instrumentos econômicos existentes. Dependem, isto sim, de verificar se a contribuição produtiva da vegetação e sua heterogeneidade territorial são relevantes para o problema estudado. A estimação proposta tampouco medirá, diretamente, efeitos sobre preços de mercado, bem-estar agregado ou lucro líquido.

\medskip
Dito isso, a análise parte das seguintes \textbf{hipóteses}:

%% H2 — Intertemporalidade e heterogeneidade espacial: investiga como o estoque e a trajetória anterior da vegetação condicionam a produtividade agrícola diante de choques climáticos, considerando diferenças territoriais e entre biomas.

\medskip
\noindent\textbf{(H2) Vegetação, clima e produtividade.}

\noindent\emph{A resposta da produtividade agrícola aos choques climáticos varia, em magnitude economicamente relevante, conforme o estoque de vegetação previamente existente no território.}

\medskip
\noindent\textbf{(H2-a) Heterogeneidade entre biomas.}

\noindent\emph{A magnitude da modificação exercida pelo estoque de vegetação sobre a resposta da produtividade agrícola aos choques climáticos é heterogênea entre biomas.}

\medskip
H2 não impõe sinal à interação entre vegetação e choque climático. A vegetação poderá estar associada ao amortecimento, à intensificação ou à ausência de alteração detectável da resposta produtiva. A interpretação do sinal dependerá também da definição do choque: déficit e excesso de precipitação serão tratados separadamente. Por sua vez, H2-a exige comparar formalmente as magnitudes entre grupos; encontrar significância em um bioma e ausência de significância em outro não demonstra que os efeitos sejam diferentes.

As hipóteses serão examinadas com dados em painel para Mato Grosso, tendo 2001--2024 como período candidato, conforme o recorte já adotado para discussão do projeto. Serão combinados rendimento de soja, estoque de vegetação e precipitação, com atenção à ordenação temporal entre condição ambiental, janela climática e colheita. A amostra final dependerá da harmonização dos municípios, da cobertura conjunta das fontes e da definição do calendário agrícola. Temperatura e déficit hídrico constituem extensões condicionadas à obtenção e ao tratamento dos dados necessários; a base atual de chuva não permite presumir sua disponibilidade.

Além desta introdução, o artigo está organizado em cinco seções. A primeira revisa a literatura sobre capital natural, vegetação, clima e produtividade. A segunda apresenta Mato Grosso como contexto empírico, descreve os dados e explicita a construção das variáveis. A terceira desenvolve a estratégia econométrica para examinar H2 e H2-a. A quarta estabelece os critérios para a futura apresentação e discussão dos resultados. Por fim, a quinta reúne as contribuições pretendidas e as limitações do estudo. Por se tratar de um projeto de qualificação, não são apresentados resultados econométricos já estimados.
